% Generated draft: requires mathematical and source review.
\documentclass[11pt]{article}
\usepackage[margin=25mm]{geometry}
\usepackage{fontspec}
\setmainfont{Latin Modern Roman}
\usepackage{amsmath,amssymb,mathtools,unicode-math}
\setmathfont{Latin Modern Math}
\usepackage{enumitem,needspace,xurl,hyperref}
\setlength{\parindent}{0pt}
\setlength{\parskip}{5pt}
\setlength{\emergencystretch}{3em}
\newcommand{\N}{\mathbb{N}}
\newcommand{\Z}{\mathbb{Z}}
\newcommand{\Q}{\mathbb{Q}}
\newcommand{\R}{\mathbb{R}}
\newcommand{\C}{\mathbb{C}}
\newcommand{\F}{\mathbb{F}}
\newcommand{\A}{\mathbb{A}}
\newcommand{\PP}{\mathbb{P}}
\newcommand{\E}{\mathbb{E}}
\newcommand{\eps}{\varepsilon}
\newcommand{\dd}{\,\mathrm{d}}
\begin{document}
\section*{UnsolvedMath ID: 10600051}
\textbf{Source problem code:} catalogue entry
\subsection{Definitions and mathematical statement}
In the usual combinatorial model, a free-link diagram is a framed four-valent graph, allowing components without vertices, where the framing specifies opposite pairs of half-edges at each vertex. Free-link equivalence is generated by the free Reidemeister moves. Let $c(D)$ be the number of vertices of a diagram $D$, and, on the usual crossing-number interpretation of minimality, put
\[
 c(L)=\min\{c(D):D\text{ represents the free link }L\}.
\]
A minimal representative has $c(D)=c(L)$.

Problem-list item 51 asks to prove or disprove the conjecture that some free link has more than one minimal representative. In this interpretation the proposed existence assertion is
\[
 \exists L\ \exists D_1,D_2:\quad
 [D_1]=[D_2]=L,\qquad
 c(D_1)=c(D_2)=c(L),\qquad D_1\not\cong D_2.
\]
Here $\cong$ must mean the intended identification of diagrams before quotienting by all free Reidemeister moves, for example framing-preserving graph isomorphism if that is the source convention. It cannot mean free-link equivalence itself, since both diagrams represent $L$. The excerpt does not specify whether additional moves among minimal diagrams, such as third Reidemeister moves, are included in the uniqueness convention, or whether oriented or labelled components are intended. These conventions must be resolved before this becomes a unique precise conjecture.
\subsection{Short English statement}
Determine whether a free link can have distinct diagrams that both realize its minimum number of crossings. The meaning of distinct minimal representatives must be fixed.
\subsection{Sources}
\begin{itemize}
\item Local registry: \url{data/unsolvedmath/problems.json}, permanent UnsolvedMath ID $10600051$, exact source code \texttt{catalogue entry}.
\item \emph{Unsolved Problems in Virtual Knot Theory and Combinatorial Knot Theory}, revised survey, problem-list item 51: \url{https://arxiv.org/abs/1409.2823}. Supplied access date: 2026-07-29; extraction: source-tex; the revised problem section contains 60 top-level items.
\item Supplied related title: ``Minimal Diagrams of Free Knots,'' with incomplete identifier \texttt{arXiv:1008.xxxx} (2010). The summary says this does not resolve the uniqueness question, but supplies no exact bibliographic reference.
\end{itemize}
\end{document}
